\chapter{Les problèmes de blocs et d'ordonnancement}
\label{chap:probleme}

\section{Description}

Une intervention constitue une tâche à positionner dans le temps et à affecter à des ressources compatibles. Elle possède notamment une durée estimée, une spécialité, un chirurgien et des besoins d'anesthésie et d'hébergement. Les ressources incluent les salles, les vacations, les équipes et les lits. La faisabilité d'un planning dépend de leurs disponibilités, des temps de préparation et de nettoyage, des priorités cliniques, des rendez-vous déjà fixés et d'une marge pour les urgences et les retards.

L'objectif ne se réduit donc pas au remplissage maximal d'une vacation. Il faut rechercher un compromis entre utilisation du bloc, stabilité de l'occupation des lits, attente des patients, respect des contraintes et robustesse face aux événements imprévus. Ces critères peuvent être contradictoires : une meilleure utilisation apparente des salles peut, par exemple, accroître les dépassements horaires ou reporter la congestion vers une unité d'aval \parencite{alamin2025,schouten2023}.

\aremplir{formaliser les ensembles, paramètres, variables de décision, contraintes et fonctions objectif.}

\section{Complexité}

\aremplir{caractériser précisément la famille du problème, sa complexité et les conséquences pour le choix des méthodes de résolution, avec des références.}

\section{État de l'art}

\subsection{Périmètre de la littérature}

La planification des blocs opératoires recouvre plusieurs décisions. À long terme, le niveau stratégique dimensionne les ressources ; à moyen terme, le niveau tactique répartit les vacations entre spécialités ou équipes ; à court terme, le niveau opérationnel sélectionne les patients, les affecte à une date et à une salle, puis détermine leur ordre de passage. La revue fondatrice de \textcite{cardoen2010} propose une classification selon le contexte du problème, les classes de patients, les critères de performance, les techniques de résolution et la prise en compte de l'incertitude. La revue proposée par \textcite{alamin2025}, fondée sur 260 publications parues entre 2000 et 2023, confirme la diversité persistante des formulations et l'absence d'une méthode dominante dans tous les contextes.

La distinction entre chirurgie ambulatoire et hospitalisation conventionnelle est structurante. Après analyse de 178 travaux, \textcite{wang2021} observent que les séjours conventionnels sont généralement associés à des interventions plus longues et plus variables ainsi qu'à davantage d'urgences, tandis que l'ambulatoire est davantage exposé aux absences des patients. Une même fonction objectif ou une même marge de sécurité ne peut donc pas être appliquée sans distinction à tous les parcours.

\subsection{Prédiction des durées}

L'apprentissage automatique intervient surtout en amont de l'optimisation pour estimer la durée d'une intervention, le risque d'annulation ou les besoins en salle de surveillance post-interventionnelle. Une revue systématique récente retient 21 études et identifie notamment les réseaux de neurones, les forêts aléatoires et XGBoost parmi les familles de modèles étudiées \parencite{merghani2025}. Ces résultats justifient la comparaison de modèles non linéaires avec une référence simple, mais ne dispensent ni d'une validation temporelle ni d'une analyse de calibration.

Le travail de \textcite{dupre2025} est particulièrement proche des cibles retenues dans ce projet : il prédit à la fois le temps d'occupation de salle et la durée de séjour à partir de données préopératoires. Sur 3\,117 interventions urgentes d'un établissement français, la forêt aléatoire obtient une erreur absolue moyenne de 32 minutes pour l'occupation de salle, et XGBoost une erreur absolue moyenne de cinq jours pour le séjour. Les auteurs qualifient toutefois les performances de modérées et signalent leur variabilité entre spécialités. Les valeurs ponctuelles doivent donc être complétées par des quantiles ou des intervalles prédictifs : ce sont les risques de dépassement de vacation et de saturation des lits, et non la seule erreur moyenne, qui intéressent l'ordonnancement.

La tendance la plus récente consiste à relier explicitement prédiction et prescription. Dans une étude sur l'arthroplastie, \textcite{lex2026} combinent une prédiction par réseau de neurones et un programme linéaire en nombres entiers. Sur des plannings simulés, leur méthode réduit la sous-utilisation de 56,2\,\% par rapport à une moyenne glissante et augmente le nombre d'interventions, mais accroît aussi les dépassements horaires de 17,2\,\%. Ce résultat illustre l'importance d'évaluer simultanément plusieurs indicateurs. Un cadre encore plus récent, disponible en ligne et rattaché à un volume de 2027, transmet au modèle d'optimisation des bornes de durée calibrées sur les résidus d'un prédicteur utilisant notamment des représentations BioBERT des codes CIM-10 ; il coordonne ensuite bloc, capacité de réveil et charge des chirurgiens \parencite{kayvanfar2027}. Cette approche prédictive--prescriptive est proche de l'ambition du présent projet, tout en laissant à traiter le séquencement intrajournalier et la validation prospective.

\subsection{Planification intégrée et incertitude}

Optimiser chaque ressource séparément peut produire une bonne solution locale mais une mauvaise performance hospitalière. La revue de \textcite{rachuba2024}, qui recense 319 travaux de recherche opérationnelle sur la planification intégrée à l'hôpital, souligne que les effets en cascade entre le bloc, les lits et les équipes ne peuvent être compris qu'en représentant conjointement le système. Elle constate aussi que les déploiements réguliers en pratique restent rares, malgré de nombreuses études de cas.

Plusieurs travaux montrent concrètement l'intérêt de déplacer les interventions électives en fonction de la capacité d'aval. \textcite{carter2024} formulent un programme en nombres entiers reliant le programme opératoire directeur et la demande attendue en lits. Sur les cas testés, le pic de demande est réduit de 3 à 19\,\% sans augmentation de capacité. À la frontière actuelle du domaine, \textcite{shehadeh2026} intègrent dans une même formulation l'affectation des interventions à une date et à une salle, leur séquencement, leurs heures de début, les lits de soins intensifs et les lits d'unité conventionnelle. Leur optimisation robuste en distribution traite conjointement l'ambiguïté sur les durées opératoires et les durées de séjour. Cette formulation constitue une référence méthodologique très proche de notre problème, mais elle suppose des capacités et des parcours postopératoires qui ne figurent pas encore dans la base fournie.

\subsection{Méthodes de résolution}

Les formulations exactes, notamment la programmation linéaire en nombres entiers, fournissent des références utiles sur les petites instances. Lorsque le nombre de patients, de salles, de périodes et de scénarios augmente, l'espace combinatoire motive l'emploi d'heuristiques, de métaheuristiques, de décompositions et de couplages avec la simulation \parencite{alamin2025}. Les travaux récents n'établissent pas de vainqueur universel : leur résultat dépend de la représentation d'une solution, du voisinage, de la fonction objectif, des contraintes et du budget de calcul.

Deux comparaisons sont directement pertinentes pour les méthodes prévues dans ce projet. \textcite{ahmadian2025} étudient un ordonnancement conjoint des salles et de la salle de surveillance post-interventionnelle, avec des durées opératoires et de réveil incertaines. Ils comparent algorithme génétique, essaim particulaire, recuit simulé et recherche taboue en utilisant un réseau de neurones comme substitut rapide à la simulation ; l'algorithme génétique obtient les meilleures solutions sur leurs instances, sans que ce classement puisse être généralisé à une autre formulation. \textcite{vieira2025} utilisent, pour leur part, un codage commun par clés aléatoires afin d'appliquer un algorithme génétique biaisé enrichi par apprentissage par renforcement, un recuit simulé et une recherche locale itérée à un problème comprenant plusieurs salles, les équipements et les disponibilités des patients et des chirurgiens. Le codage partagé est une idée importante pour notre protocole : comparer plusieurs métaheuristiques sur la même représentation et avec le même évaluateur limite les biais expérimentaux.

\subsection{Positionnement du projet}

Le projet s'inscrit donc dans la famille des approches \emph{prédictives--prescriptives} et de planification intégrée : prédire, avant l'admission, le temps d'occupation de salle et la durée de séjour, puis utiliser ces estimations pour proposer des dates compatibles avec les vacations et pour lisser la demande en lits. Son intérêt ne réside pas dans l'affirmation qu'une métaheuristique serait toujours supérieure, mais dans une comparaison reproductible de la recherche taboue, du recuit simulé et de l'algorithme génétique sur une formulation commune, avec des critères opérationnels et une analyse de robustesse.

Cette ambition doit rester proportionnée aux données disponibles. La base historique permet d'apprendre les deux durées cibles, mais elle ne renseigne ni l'identifiant des salles, ni les vacations planifiées, ni la capacité quotidienne des lits et de la salle de réveil, ni les annulations et arrivées urgentes. Ces éléments devront être obtenus auprès de l'établissement ou introduits comme hypothèses explicites dans des instances simulées. L'évaluation comparera au minimum les modèles prédictifs à une moyenne ou une médiane historique, et les métaheuristiques à une règle gloutonne ainsi qu'à une résolution exacte sur les petites instances. Les indicateurs couvriront l'utilisation et les dépassements de vacation, l'attente ou le report des patients, la variabilité et les pics d'occupation des lits, les violations de contraintes, le temps de calcul et la sensibilité aux aléas. Ce choix répond à la recommandation de \textcite{schouten2023} de ne pas améliorer un indicateur isolé au détriment de la qualité globale du système.
